\section{Introduction}

A circle domain is a nonempty connected open subset of the Riemann
sphere $\Chat=\C\cup\{\infty\}$ whose complementary components are
points or closed round disks; the complement may be empty.  Circle
boundaries give a particularly concrete conformal model of a plane
domain.  Their usefulness depends on two separate questions: whether
such a model exists, and whether its conformal structure determines
its placement in the sphere up to a M\"obius transformation.  This paper
addresses the second question under conformal removability of the
boundary.

A compact set $E\subset\Chat$ is \emph{conformally removable} if every
orientation-preserving homeomorphism of the sphere that is conformal
on $\Chat\setminus E$ is M\"obius.  A circle domain is
\emph{conformally rigid} if every conformal equivalence from it onto
another circle domain is the restriction of a M\"obius transformation.
The definitions have different starting data: removability concerns a
map already defined on the sphere, whereas rigidity concerns only a
map between the domains.

\begin{theorem}\label{thm:main}
Every circle domain with conformally removable boundary is conformally
rigid.  More explicitly, if $\Omega,\Omega'\subset\Chat$ are circle
domains, $\partial\Omega$ is conformally removable, and
$f\colon\Omega\to\Omega'$ is a conformal equivalence, then there is a
M\"obius transformation $M$ such that $f=M|_\Omega$.
\end{theorem}

The theorem allows arbitrary cardinality of the complementary component
set.  The required boundary control is a conclusion of the argument;
no boundary extension of $f$ is part of the hypothesis.

\subsection{Rigidity and removability}

He and Schramm proved circle-domain uniformization for countably
connected domains and rigidity for countably connected circle
domains~\cite{HeSchramm1993}.  Their subsequent rigidity theorem permits
uncountably many complementary components when the boundary has
$\sigma$-finite linear measure~\cite{HeSchramm1994}.  They conjectured
that a circle domain is rigid if and only if its boundary is conformally
removable.  The forward implication in this proposed equivalence is
false: Rajala constructed a rigid circle domain with a non-removable
boundary~\cite[Theorem~1.1]{Rajala2025}.  Theorem~\ref{thm:main}
establishes the remaining removability-to-rigidity implication
positively.

Several intermediate results clarify the role of boundary geometry.
Younsi studied the formulations involving removability of the whole
boundary and of its Cantor subsets, and proved that conformal rigidity
is equivalent to quasiconformal rigidity~\cite[Theorems~3--5]{Younsi2016}.
Ntalampekos and Younsi proved rigidity under square integrability of the
quasihyperbolic distance in a bounded portion of the domain; this
includes H\"older and John circle domains~\cite[Theorem~1.6 and
Corollary~1.7]{NY2020}.  Ntalampekos subsequently proved rigidity when
every compact subset of the set of point boundary components is
\emph{countably negligible for extremal distance}, or CNED
\cite[Theorem~1.2]{Ntalampekos2023}.  For every pair of disjoint nondegenerate continua, the CNED condition
says that restricting the joining curves to those that meet the
exceptional set in at most countably many points leaves their conformal
modulus unchanged.  These results develop
methods for controlling boundary traces under explicit geometric or
extremal-distance hypotheses.

Continuous Sobolev extension is a second relevant line of work.
Jones and Smirnov related removability for continuous Sobolev functions
to quasiconformal removability and proved criteria based on the geometry
near a domain boundary~\cite[Definition~2 and Theorems~2,~4]{JonesSmirnov2000}.
For closed CNED sets, Ntalampekos proved continuous $W^{1,n}$ extension
in dimension $n$~\cite[Theorem~1.4]{NtalampekosCNED2024}.
Our intermediate Theorem~\ref{thm:dust-sobolev} treats compact totally
disconnected planar sets under conformal removability itself: a scalar
function continuous across such a set, with finite Dirichlet energy
on its complement, is globally Sobolev.  Continuity across the set is
an essential hypothesis of this intermediate statement.

\subsection{The proof and its quantitative input}

The first difficulty is to use removability before a conformal map is
known to extend to the sphere.  We handle it initially for the complement
of a compact totally disconnected set $B$.  A normalized conformal map
$h$ of $\Chat\setminus B$ may replace a point of $B$ by a nonpoint
complementary component in its image.  Component correspondence supplies
a continuous inverse collapse map $\pi_h$ on the target sphere.  For a
continuous scalar function $u$ on the source sphere, the function
$u\circ\pi_h$ therefore has the prescribed constant value on each target
component.  We call $h$ a model retaining $u$ when this continuation is
globally Sobolev and its gradient vanishes on the omitted set.

The companion paper \cite{KoebeCompanion} proves a quantitative
retained-test transfer and selection theorem.  We use its scalar form
in Section~\ref{sec:models}: compatible finite lists of scalar functions
are retained exactly on exhausting cores, while the remaining energy
and the square sum of boundary oscillations tend to zero.  Compatibility
here is the condition that the joint image of $B$ under each pair of
tests has planar area zero.  Filling the finite round holes and passing
to a weak limit produces the required models.  A separate contour
argument identifies their values on all omitted components, including
point components whose union can have positive area.  This is the
quantitative part of the uniformization theorem needed in this paper.

The next step varies the model.  Among maps retaining a fixed finite
list, we nearly maximize the real part of the first Laurent coefficient
at infinity.  A conductivity deformation, constructed through the
classical Beltrami equation, then gives a scalar height close to the
ordinary vertical coordinate, with arbitrarily small Dirichlet error.
A level-set argument makes this height compatible with the retained
tests.  After flattening the countably many values belonging to nonpoint
holes, a horizontal coordinate can be added as well.  Together the two
coordinates yield barriers around all source ends with arbitrarily
small energy.

We save these barriers while enlarging the finite test lists.  If a
limiting model had a nonpoint omitted component, its projection on one
axis would contain an interval of positive length.  Every transverse
line over that interval would force a saved barrier to change from one
to zero in a bounded distance.  The resulting fixed lower bound on its
energy contradicts the construction.  Thus every omitted component is
a point, and component correspondence extends the limiting model to a
sphere homeomorphism.  Removability makes the normalized limit the
identity; weak Sobolev compactness then proves the scalar extension
theorem.

Finally we return to a conformal equivalence $f:\Omega\to\Omega'$.
Point components may correspond to disks, and disks to points, so direct
continuation of a coordinate of $f$ need not be continuous.  We mask the
countably many disk-containing pairs and apply the scalar theorem to
the remaining common point components.  The masks show that their
traces have null axial projections on almost every line.  This use of
linewise variation is related to the Sobolev-removability arguments of
Jones and Smirnov~\cite{JonesSmirnov2000} and to the boundary-control
methods in~\cite{NY2020,Ntalampekos2023}; the mask construction here
accommodates all common point components simultaneously.  Integrating
the resulting slice estimates gives a flux inequality.  Comparison with
source and target area bounds forces equality, and hence forces $f$ to
be M\"obius.

Section~\ref{sec:preliminaries} constructs component correspondence and
records normalization bounds.  Section~\ref{sec:models} obtains models
retaining prescribed tests.  Section~\ref{sec:conductivity} supplies
the conductivity and level-set estimates, and
Section~\ref{sec:sobolev} uses them to prove scalar Sobolev extension.
Section~\ref{sec:traces} constructs the masks and proves the trace
statement.  Section~\ref{sec:rigidity} completes the flux and area
argument.
